\section{The autonomous type-2 probability}
\label{sec:autonomous}

The second equation in \eqref{eq:backward-scaled} does not depend on $\widehat S$: it is autonomous. This is the small structural fact on which the full calculation rests. We solve it first: factor the right-hand side, introduce a decaying coordinate $z_2$, and write $\widehat G$ as a rational function of $z_2$. That rational time dependence will then drive the type-1 equation.

Constant solutions of \cref{eq:G-scaled} satisfy a quadratic equation. Setting
\begin{equation}
\label{eq:Delta2}
\hat\Delta_2=
\sqrt{(\hat\lambda_2+\hat\mu_2+\hat\delta_2)^2
-4\hat\lambda_2\hat\mu_2},
\end{equation}
the two roots are
\begin{equation}
\label{eq:type2-roots}
r_{2,\pm}
=\frac{\hat\lambda_2+\hat\mu_2+\hat\delta_2
\pm\hat\Delta_2}{2\hat\lambda_2}.
\end{equation}
They satisfy
\[
r_{2,-}+r_{2,+}
=\frac{\hat\lambda_2+\hat\mu_2+\hat\delta_2}{\hat\lambda_2},
\qquad
r_{2,-}r_{2,+}=\frac{\hat\mu_2}{\hat\lambda_2},
\]
and therefore the right-hand side factors as
\begin{equation}
\label{eq:G-factored}
\widehat G'
=\hat\lambda_2(\widehat G-r_{2,-})(\widehat G-r_{2,+}).
\end{equation}

There is also a probabilistic ordering among these equilibria. The quadratic is non-negative at zero and equals $-\hat\delta_2<0$ at one, while its leading coefficient is positive. Therefore
\begin{equation*}
0\leq r_{2,-}<1<r_{2,+}.
\end{equation*}
Note that a long-time probability must lie in $[0,1]$, so only the lower root $r_{2,-}$ can be an admissible attractor; the upper root lies outside that interval.
\Cref{fig:type2-phase} makes this geometry and the subsequent relaxation explicit.

\begin{figure}[htbp]
  \centering
  \includegraphics[width=\textwidth]{figures/fig06_type2_phase/fig06.pdf}
  \caption{Geometry of the autonomous type-2 probability.
  (a)~Right-hand side $f(G)=\lambda_2(G-r_{2,-})(G-r_{2,+})$, with $0\le r_{2,-}<1<r_{2,+}$ so only $r_{2,-}$ is an admissible long-time probability.
  (b)~Phase line: $r_{2,-}$ attracting, $r_{2,+}$ repelling; $G(0)=1$ is drawn into $r_{2,-}$.
  (c)~Relaxation of $G(t)$ from $1$ toward $r_{2,-}$.
  Parameters: $\lambda_2=0.9$, $\mu_2=0.35$, $\delta_2=0.20$.}
  \label{fig:type2-phase}
\end{figure}

We shall use the root gap and its associated decay rate repeatedly. Define
\begin{equation}
\label{eq:type2-alpha-q}
\alpha_2=r_{2,+}-r_{2,-}
=\frac{\hat\Delta_2}{\hat\lambda_2},
\qquad
\hat q_2=-\hat\lambda_2\alpha_2=-\hat\Delta_2.
\end{equation}
Thus $\alpha_2>0$ and $\hat q_2<0$.

\begin{proposition}[Exact type-2 probability]
\label{prop:G-solution}
Let
\begin{equation}
\label{eq:type2-z}
z_{2,0}=\frac{r_{2,-}-1}{r_{2,+}-1}<0,
\qquad
z_2(\hat t)=z_{2,0}\e^{\hat q_2\hat t}.
\end{equation}
Then the solution of \cref{eq:G-scaled} with $\widehat G(0)=1$ is
\begin{equation}
\label{eq:G-solution}
\widehat G(\hat t)
=r_{2,-}-\alpha_2\frac{z_2(\hat t)}{1-z_2(\hat t)}.
\end{equation}
In physical time, with $\hat t=\lambda_1t$ and the scaled rates of \cref{eq:scaled-rates},
\begin{equation*}
G(t)
=r_{2,-}-\alpha_2\frac{z_2(\lambda_1t)}{1-z_2(\lambda_1t)}.
\end{equation*}
\end{proposition}

\begin{proof}
Separate \cref{eq:G-factored} and use partial fractions:
\[
\frac{\dd\widehat G}
{(\widehat G-r_{2,-})(\widehat G-r_{2,+})}
=\hat\lambda_2\,\dd\hat t.
\]
After integration,
\[
\log\!\left(
\frac{\widehat G-r_{2,-}}{\widehat G-r_{2,+}}
\right)
=-\hat\lambda_2\alpha_2\hat t+\text{constant}.
\]
The left-hand ratio is the natural coordinate produced by separation of variables. The initial condition makes it equal to $z_{2,0}$ at zero, so it equals $z_2(\hat t)$ at time $\hat t$. Solving the resulting linear fractional equation for $\widehat G$ gives \cref{eq:G-solution}.
\end{proof}

Thus $z_2$ is simply that separated ratio, started from $\widehat G(0)=1$ and decaying exponentially toward zero as $\widehat G$ approaches $r_{2,-}$. Three immediate checks are
\begin{equation*}
\widehat G(0)=1,
\qquad
\widehat G'(0)=-\hat\delta_2,
\qquad
\widehat G(\hat t)\longrightarrow r_{2,-}
\quad(\hat t\to\infty).
\end{equation*}
The last limit has a direct meaning. Starting from one type-2 individual, $r_{2,-}$ is the probability that catastrophe is avoided forever. It can be positive because the lineage may die out ordinarily before it triggers anything.

An expanded form may be useful for direct implementation:
\begin{equation*}
\widehat G(\hat t)=
\frac{r_{2,-}(r_{2,+}-1)
+r_{2,+}(1-r_{2,-})
\e^{-\hat\lambda_2(r_{2,+}-r_{2,-})\hat t}}
{(r_{2,+}-1)
+(1-r_{2,-})
\e^{-\hat\lambda_2(r_{2,+}-r_{2,-})\hat t}}.
\end{equation*}
In physical time the same expression is $G(t)=\widehat G(\lambda_1t)$. The compact coordinate $z_2$ is the better choice for the driven equation: it is always negative, moves monotonically toward zero, and carries all of the time dependence of $\widehat G$ in a rational way (\cref{fig:G-and-z2}). With $\widehat G$ known, the type-1 equation is now a forced Riccati equation in $\widehat S$. In the next section we substitute this $G$ and use $z_2$ as the independent variable.

\FloatBarrier

\begin{figure}[htbp]
  \centering
  \includegraphics[width=\textwidth]{figures/fig07_G_and_z2/fig07.pdf}
  \caption{(a)~$G(t)$ from $1$ toward $r_{2,-}$.
  (b)~Coordinate $z_2(\lambda_1t)$ increasing from $z_{2,0}<0$ toward $0$.
  (c)~Linear-fractional map $G=r_{2,-}-\alpha_2 z_2/(1-z_2)$.
  Parameters: $\lambda_1=1$, $\lambda_2=0.85$, $\mu_2=0.40$, $\delta_2=0.18$.}
  \label{fig:G-and-z2}
\end{figure}

